\section{\texorpdfstring{Decision procedures: \href{https://lean-lang.org/doc/reference/latest/Tactic-Proofs/Tactic-Reference/}{\texttt{decide}}, \href{https://lean-lang.org/doc/reference/latest/Tactic-Proofs/Tactic-Reference/}{\texttt{omega}}, \href{https://lean-lang.org/doc/reference/latest/Tactic-Proofs/Tactic-Reference/}{\texttt{norm\_num}}}{Decision procedures: decide, omega, norm\_num}}\label{sec:12-working-efficiently:02-decision-procedures:1}

For goals that are \emph{decidable} --- where ``true or false'' can be settled by a terminating algorithm instead of a hand-built argument --- Lean has tactics that just run that algorithm:

\begin{itemize}
\tightlist
\item
  \textbf{\texttt{decide}} --- evaluates a \texttt{Decidable} proposition to \texttt{true}/\texttt{false} directly. It works well for small, closed (no free variables) propositions, for example \texttt{(7 : Nat) ∣ 21} or \texttt{¬ (3 = 5)}. \texttt{decide} should not be used on propositions with free variables or unbounded search: it can time out, or worse, produce a correct but useless proof term that reveals nothing.
\item
  \textbf{\texttt{omega}} --- a decision procedure for \emph{linear} arithmetic over \texttt{Nat}/\texttt{Int} (goals built from \texttt{+}, subtraction, \texttt{≤}, \texttt{<}, \texttt{=}, and multiplication by \emph{literal constants} only --- \texttt{omega} handles \texttt{3 * n} fine, but not \texttt{n * m} for two unknown variables \texttt{n}, \texttt{m}; multiplying two unknown variables together falls outside what it can decide). For a goal that is ``some linear inequality or equality about integers,'' \texttt{omega} should be reached for before deriving the fact by hand. This is exactly the kind of fact a decision procedure handles better than a custom \href{https://lean-lang.org/doc/reference/latest/Tactic-Proofs/Tactic-Reference/}{\texttt{rw}} chain, and a hand-derived proof teaches nothing that \texttt{omega}'s existence does not already establish.
\item
  \textbf{\texttt{norm\textbackslash{}\_num}} --- normalizes and evaluates numerical expressions (arithmetic on concrete numerals, including some with \texttt{+}, \texttt{*}, \texttt{\textbackslash{}\textasciicircum{}}, \texttt{≤} on literals).
\end{itemize}

Here is the judgment call: \textbf{use a decision procedure whenever the goal falls entirely inside its decidable fragment; use explicit \texttt{rw}/\texttt{have} reasoning whenever the goal is about a \emph{general}, unspecified structure} (like an arbitrary \texttt{Group G} or \texttt{Ring R}). No decision procedure applies there, because there is no concrete computation to run, only axioms to combine. Chapters 7 and 9's group/ring theorems are all of this second kind, which is exactly why they needed hand-built proofs instead of \texttt{omega}.

\begin{mathreading}

A proposition \(P\) is \texttt{Decidable} when there is an algorithm that computes its truth value. In categorical terms, this is a map \(P \to \{\top, \bot\}\) whose two preimages are ``proof of \(P\)'' and ``proof of \(\neg P\).'' So \texttt{decide} is the constructive statement ``\(P \vee
\neg P\) holds \emph{and} we can tell which one,'' available exactly on the decidable fragment (closed numeric claims like \(7 \mid 21\)). \texttt{omega} decides \emph{Presburger arithmetic}, the first-order theory of \((\mathbb{Z},
+, <)\), which is famously decidable, and \texttt{norm\textbackslash{}\_num} evaluates concrete numerals. These apply only to statements with no free structure left to fill in. A theorem about an arbitrary group has no finite truth table to compute, and thus must be \emph{proved} from the axioms instead of \emph{decided}.

\end{mathreading}

\begin{quote}
Read more: the Lean/Mathlib documentation for \texttt{decide}, \texttt{omega}, and \texttt{norm\textbackslash{}\_num} (searchable in the Mathlib docs or via \texttt{\textbackslash{}\#help tactic omega} inside a Lean file) covers each tactic's exact decidable fragment and performance in more depth than this book needs.
\end{quote}
